\documentclass[12pt]{article}
\input{../Paquetes/paquetes.tex}

\newcommand{\Titulo}{1er Lista de problemas y ejercicios}

\begin{document}
	\input{../Portada/PORTADA_TC.tex}
	%
    \pagenumbering{Roman}
    %
    \newpage
    \thispagestyle{empty}
    \setcounter{page}{1}
    %\section*{Resumen}
    % ============================================================
% Resumen (abstract)
% ============================================================

\begin{abstract}
    \noindent En esta lista de problemas se estudian los fundamentos de las cadenas y los
    lenguajes formales sobre un alfabeto $\Sigma$. Se trabajó con los conjuntos
    $\Sigma^n$, $\Sigma^*$ y $\Sigma^+$, la longitud y la concatenación de cadenas, y
    operaciones entre lenguajes como la unión, la intersección, las potencias, la
    reflexión y las clausuras de Kleene y positiva. Mediante el Principio de
    Inducción Matemática se demostraron resultados como $|\Sigma^n| = |\Sigma|^n$,
    $|u^n| = n|u|$ y $(uv)^R = v^R u^R$, además de las propiedades de $A^*$ y $A^+$.
    Entre las observaciones principales destacan que la concatenación es asociativa
    pero no conmutativa, aunque $|uv| = |vu|$; que en general $L^* \neq L$; y que el
    lenguaje vacío $\varnothing$ debe distinguirse de $\{\lambda\}$. Estos resultados
    constituyen la base para el estudio posterior de los autómatas finitos, las
    expresiones regulares y las gramáticas formales.
\end{abstract}
    \addcontentsline{toc}{section}{Resumen}
    %
    \newpage
    %\thispagestyle{empty}
    \tableofcontents
    %%
    %\newpage
    %\thispagestyle{empty}
    %\listoffigures
    %\addcontentsline{toc}{section}{Índice de figuras}
    %%
    %\newpage
    %\thispagestyle{empty}
    %\listoftables
    %\addcontentsline{toc}{section}{Índice de tablas}
    %
    \newpage
    \thispagestyle{plain}
    \pagenumbering{arabic}
    \setcounter{page}{1}
    \section*{Introducción}
    % ============================================================
% Introducción
% ============================================================

La Teoría de la Computación estudia qué problemas pueden resolverse de forma
mecánica y con qué recursos. Para plantear estas preguntas con precisión es
necesario contar primero con un lenguaje matemático que permita describir la
información que procesa una máquina. Ese lenguaje se construye a partir de tres
nociones elementales: el \textbf{alfabeto} $\Sigma$, un conjunto finito y no vacío de
símbolos; la \textbf{cadena} o palabra, una secuencia finita de símbolos de
$\Sigma$; y el \textbf{lenguaje}, cualquier conjunto de cadenas, es decir, cualquier
subconjunto de $\Sigma^*$.

La presente lista de problemas y ejercicios tiene como objetivo afianzar estas
nociones y las operaciones que se definen sobre ellas, las cuales constituyen la
base sobre la que más adelante se estudiarán los autómatas finitos, las
expresiones regulares y las gramáticas formales. En particular, se trabaja con:

\begin{itemize}
    \item \textbf{Estructuras ordenadas:} los conjuntos $\Sigma^n$, $\Sigma^*$ y
    $\Sigma^+$, la longitud y la concatenación de cadenas, y la pertenencia de
    cadenas a lenguajes definidos por comprensión, junto con sus uniones e
    intersecciones.
    \item \textbf{Clausura de Kleene y clausura positiva:} demostración de las
    propiedades de $A^*$ y $A^+$, como $A^+ = A^* A$, $(A^*)^* = A^*$ y
    $(A \cup B)^* = (A^* B^*)^*$.
    \item \textbf{Reflexión de un lenguaje:} la cadena inversa $u^R$, el lenguaje
    inverso $A^R$ y propiedades como $(A \cdot B)^R = B^R A^R$ y
    $(A^*)^R = (A^R)^*$.
    \item \textbf{Longitud de cadenas y $n$-cadenas:} resultados de conteo como
    $|w^{n+m}| = |w^n| + |w^m|$ y $|\Sigma^n| = |\Sigma|^n$.
    \item \textbf{Prefijos, sufijos y potencias de un lenguaje:} descomposición de
    una cadena y definición recursiva de $A^n$.
    \item \textbf{Ejercicios de aplicación:} cálculo de concatenaciones, potencias,
    cerraduras y reflexiones de lenguajes concretos, así como la
    determinación del alfabeto sobre el que está definido un lenguaje.
\end{itemize}

La herramienta principal a lo largo del trabajo es la \textbf{demostración formal},
en especial el \textbf{Principio de Inducción Matemática}, que resulta natural en
este contexto porque las cadenas, su longitud, su concatenación y su reverso se
definen de manera recursiva. Cada ejercicio incluye su enunciado y el
procedimiento o razonamiento que justifica la respuesta.
    %
	\newpage
    \section{Problemas de Estructuras Ordenadas}
        Sea $\Sigma = {0,1}$. En los problemas 1-4, encuentre una cadena $w \in A$ con las características indicadas, cuando sea posible.

\begin{enumerate}
    \item $A = \Sigma^{*}, w : |w| = 5$
        \subitem $R: $ Siendo que la cerradura de Kleene comprende todas las combinaciones posibles de cadenas en el alfabeto $(\Sigma)$, en este caso {$0,1$}. Tenemos que una longitud de la cadena $w$ que sea igual a $5$ puede ser vista como: $w = abcde$ con $a,b,c,d,e \in \Sigma$. Por lo tanto, también esto puede ser descrito como $w \in \Sigma^{5}$ contanto con $2^5 = 32$ cadenas posibles. Cadena tomada al azar: $10110$
    %
    \item $A = \Sigma^{+}, w : |w| = 0$
        \subitem $R: $ Siendo que la cerradura positiva de un alfabeto excluye la cadena vacía; por tanto, este problema no tendría solución posible, ya que la indicación no es compatible con sus límites. $No$ $solution.$
    %    
    \item $A = \Sigma^{*}, w : |w| = 0$
        \subitem $R: $ Con la notación de cerradura de Kleene podemos acceder a la cadena vacío $\lambda$ que permite la realización de este ejercicio. Donde si $w = \lambda \rightarrow |w| = 0$.
    %
    \item $A = \Sigma^{5}, w : w_0 = 0$
        \subitem $R :$ Teniendo en cuenta la notación de un alfabeto restringido a 5 caracteres por cadena posible cuyo primer dígito sea 0, podemos proponer cualquier cadena tal que $w = abcde \in \Sigma^5, a = 0$. Y podemos proponer de forma aleatoria la cadena $00101$. Teniendo $2^4 = 16$ cadenas posibles, ya que el primer caracter estará fijado en 0.
\end{enumerate}

Sea $\Sigma = \{a,b\}; L_1 = \{w\in\Sigma^* : w_0 = a\}; L_2 = \{w\in\Sigma^* : w_n = b\}, n = |w|$. Dadas las cadenas $w=abbbaa; \pi = aaabb; \tau = baaba.$, realice los siguientes ejercicios.

\begin{enumerate}
    \setcounter{enumi}{4}
    \item Determine $|\pi|$.
        \subitem $R: $ La longitud de la cadena $\pi$ es igual a $5$, por ende, pertenece a $\Sigma^5$.
    %
    \item Encuentre $w\pi$.
        \subitem $R: $ La concatenacion entre $w$ y $\pi$ en ese orden es $abbbaa aaabb$
    %
    \item Encuentre $\pi w$.
        \subitem $R: $ La concatenacion entre $\pi$ y $w$ en ese orden es $aaabb abbbaa$
    %
    \newpage
    \item Verifique si $(w\pi)\tau = w(\pi\tau)$.
        \subitem $R: $ Demostrando la asociatividad de la concatenación de cadenas, tenemos como primera cadena:
        \begin{center}
            $(w\pi)\tau = (abbbaa aaabb) baaba = abbbaaaaabbbaaba$\\
        \end{center}
        Para el desarrollo del lado derecho de la igualdad:
        \begin{center}
            $w(\pi\tau) = abbbaa (aaabb baaba) = abbbaaaaabbbaaba$.\\
        \end{center}
        Por lo tanto se verifica la igualdad de ambas concatenaciones.\\
    %
    \item Proporcione una palabra en $\Sigma^7$.
        \subitem $R: $ Se propone $abbabba$, que satisface la restricción de combinación de 7 índices.
    %
    \item Proporcione una palabra en $L_1 \cap \Sigma^5$.
        \subitem $R: $ Se propone $ababa$ que satisface la restricción de $L_1$ (cuya primera letra debe ser $a$) y la restricción de $\Sigma^5$ (combinación de 5 índices).
    %
    \item Proporcione una palabra en $L_2 \cap \Sigma^4$.
        \subitem $R: $ Se propone $babb$ que satisface la restricción de $L_2$ (cuya última letra debe ser $a$) y la restricción de $\Sigma^4$ (combinación de 5 índices). y la restricción de $\Sigma^4$ (combinación de 4 índices).
    %
    \item Proporcione una palabra en $L_1 \cap L_2$.
        \subitem $R: $ Se propone $ab$ que satisface la restricción de $L_1$ (cuya primera letra debe ser $a$) y la restricción de $L_2$ (cuya última letra debe ser $b$).
    %
    \item Determine $L_1 \cap \Sigma^0$
        \subitem $R: $ Se determina como $\emptyset$ la intersección entre $L_1$ y $\Sigma^0$, ya que $\Sigma^0 = \{\lambda\}$ y $L_1$ restringe a la existencia de algún caracter $\"a\"$ en su primer índice. Por lo tanto la operación no es posible. $No$ $Solution$.
\end{enumerate}

Sea $\Sigma = \{a,b,c,d\}$ y $w=abbca; \pi=abb; \tau=ca$. Consideremos los siguientes lenguajes:

$L_1 = \{\rho \in \Sigma^* : |\rho| \leq 5\}$.\\
$L_2 = \{a^n b^m : n < m\}$.\\
$L_3 = \{\rho \in \Sigma^* : \rho$ tiene más letras $b$ que $a$ $\}$.\\

Indique si las siguientes afirmaciones (ejercicios 14-46) son verdaderas o falsas:

\begin{enumerate}
    \setcounter{enumi}{13}
    \item $|w| = 3$
        \subitem $R: $ $Falso$
    %
    \item $|w| = 5$
    \subitem $R: $ $Verdadero$
    %
    \item $w = \tau\pi$
    \subitem $R: $ $Falso$
    %
    \item $w = \pi\tau$
    \subitem $R: $ $Falso$
    %
    \item $w(\pi\tau) = (w\pi)\tau$
    \subitem $R: $ $Verdadero$
    %
    \item $w(\tau\pi) = (w\pi)\tau$
    \subitem $R: $ $Verdadero$
    %
    \item $\lambda \in \Sigma^+$
    \subitem $R: $ $Falso$
    %
    \item $\lambda \in \Sigma^*$
    \subitem $R: $ $Verdadero$
    %
    \item $w \in \Sigma^5$
    \subitem $R: $ $Verdadero$
    %
    \item $w \in \Sigma^3$
    \subitem $R: $ $Falso$
    %
    \item $\pi\tau \in \Sigma^5$
    \subitem $R: $ $Verdadero$
    %
    \item $\pi\lambda \in \Sigma^3$
    \subitem $R: $ $Verdadero$
    %
    \item $a^8 w \in \Sigma^8$
    \subitem $R: $ $Falso$
    %
    \item $a^3 w \in \Sigma^8$
    \subitem $R: $ $Verdadero$
    %
    \item $a^3 c^3 \in \Sigma^9$
    \subitem $R: $ $Falso$
    %
    \item $|a○^2 b^3| = |c^6|$
    \subitem $R: $ $Falso$
    %
    \item $|a^4 b^3 c| = |a^3 b^2 c^3|$
    \subitem $R: $ $Falso$
    %
    \item $ab \in L_1$
    \subitem $R: $ $Verdadero$
    %
    \item $ab \in L_2$
    \subitem $R: $ $Falso$
    %
    \item $ab \in L_3$
    \subitem $R: $ $Falso$
    %
    \item $abab^2 \in L_1$
    \subitem $R: $ $Verdadero$
    %
    \item $abab^2 \in L_2$
    \subitem $R: $ $Verdadero$
    %
    \item $abab^2 \in L_3$
    \subitem $R: $ $Verdadero$
    %
    \item $abcad \in L_1$
    \subitem $R: $ $Verdadero$
    %
    \item $abcad \in L_2$
    \subitem $R: $ $Falso$
    %
    \item $abcad \in L_3$
    \subitem $R: $ $Falso$
    %
    \item $L_1 \subseteq L_2$
    \subitem $R: $ $Falso$
    %
    \item $L_1 \subseteq L_3$
    \subitem $R: $ $Falso$
    %
    \item $L_2 \subseteq L_3$
    \subitem $R: $ $Verdadero$
    %
    \item $L_3 \subseteq L_2$
    \subitem $R: $ $Verdadero$
    %
    \item Si $\alpha \in \Sigma^*$, entonces $\alpha\lambda = \lambda\alpha$
    \subitem $R: $ $Verdadero$
    %
    \item Si $\alpha \in \Sigma^3$ y $\beta \in \Sigma^4$, entonces $|\alpha\beta| = 7$
    \subitem $R: $ $Verdadero$
    %
    \item Si $\alpha \in \Sigma^3$ y $\beta \in \Sigma^4$, entonces $|\alpha\beta| = 7$
    \subitem $R: $ $Falso$
\end{enumerate}

En caso de que sea posible, encuentre cinco palabtas de cada uno de los lenguajea que se definen en los problemas 4-56, usando el alfabeto $\Sigma = \{0,1,2,3,4,5,6\}$

\begin{enumerate}
    \setcounter{enumi}{46}
    \item $L_1 = \{01^n : n \in \mathbb{N}\}$
        \subitem $R: $ $0\underbrace{11\ldots1}_{n-\text{veces}}$
    %
    \item $L_2 = \{0^n 1^n 2^n : n \in \mathbb{N}\}$
        \subitem $R: $ 
            $\underbrace{0\ldots0}_{n-\text{veces}}$
            $\underbrace{1\ldots1}_{n-\text{veces}}$
            $\underbrace{2\ldots2}_{n-\text{veces}}$
    %
    \item $L_3 = \{w \in \Sigma^* : |w| = 6\}$
        \subitem $R: $ $012345$ (cadena tomada al azar)
    %
    \item $L_4 = \{w \in \Sigma^* :$ la suma de los digitos de $w$ es igual a $6\}$
        \subitem $R: $ $012012$. Donde $0+1+2+0+1+2 = 6$.
    %
    \item Determine $L_1 \cap L_2$.
        \subitem $R: \emptyset$ No es posible ya que $L_2$ necesita el uso de la letra $2$, mientras que $L_1$ requiere un $0$ al inicio.
        
        \begin{center}
            Let x: $0^0 1^0 2^0$; Let y: $01^0$\\
            $x = \lambda$, mientras $y = 0$ \\
            $\therefore$ $x \neq y$
        \end{center}
    %
    \item Determine $L_1 \cap L_3$.
        \subitem $R: $ $w = 011111$. Satisface las reglas de ambos lenguajes, donde $w_0 = 0$ y donde $|w| = 6$
    %
    \item Determine $L_1 \cap L_4$.
        \subitem $R: $ $w = 0111111$. Satisface las reglas de ambos lenguajes, donde $w_0 = 0$ y donde $0+1+1+1+1+1+1 = 6$.
    %
    \item Determine $L_2 \cap L_3$.
        \subitem $R: $ $w = 001122$. Satisface las reglas de ambos lenguajes, donde $\{0^a 1^b 2^c : a = b = c\}$ y con $|w| = 6$.
    %
    \item Determine $L_2 \cap L_4$.
        \subitem $R: $ $w = 001122$. Satisface las reglas de ambos lenguajes, donde $0+0+1+1+2+2 = 6$ y con $\{0^a 1^b 2^c : a = b = c\}$
    %
    \item Determine $L_3 \cap L_4$.
        \subitem $R: $ $w = 000501$. Satisface las reglas de ambos lenguajes, donde $0+0+0+5+0+1 = 6$ y donde $|w| = 6$.
\end{enumerate}
    \newpage
    \section{Propiedades de la Clausura de Kleene y Positiva}
        Sea $A$ un lenguaje sobre $\Sigma$, es decir, $A \subseteq \Sigma{*}$. Las siguientes son las propiedades de la clausura de Kleene $(A^*)$ y la clausura positiva $(A^+)$. Demuestre las 9 propiedades (57-65).

\begin{enumerate}
    \setcounter{enumi}{56}
    \item $A^+ = A^* * A = A * A^*$.
        \subitem Tenemos que $A^+$ contiene todas las cadenas positivas, es decir, todas las cadenas de $A^*$ excepto $\lambda$. Si concatenamos $A^*$ --que sí contiene $\lambda$-- con $A$, se agrega al menos un elemento de $A$.
        
        Entonces, tomando como ejemplo $A=\{a\}$:
        \begin{center}
            $A^* = \{\lambda, a, aa, aaa, \ldots\}$
        \end{center}
        
        Y tenemos que:
        \begin{center}
            $A^+ = \{a, aa, aaa, \ldots\}$
        \end{center}
        
        Por tanto, al concatenar $A^*$ con $A$, se agrega al menos un elemento, por lo que no es posible obtener la cadena vacía $\lambda$:
        \begin{center}
            $\therefore A^*A=A^+$
        \end{center}
        
        De manera análoga, al concatenar $A$ con $A^*$ también se obtiene una cadena de longitud positiva:
        \begin{center}
            $\therefore AA^*=A^+$
        \end{center}
        
        Por lo tanto:
        \begin{center}
            $\boxed{A^+=A^*A=AA^*}$
        \end{center}
    %
    \item $A^* * A^* = A^*$.
        \subitem Sea $A = \{a\}$ y $w$ es una cadena de $A$.\\
        
        Tenemos:
        \begin{center}
            $A^* = \{\lambda, a, aa, aaa, \ldots\}$
        \end{center}
        
        Y recordemos también que:
        \begin{center}
            $A^* = A^0 \cup A^1 \cup A^2 \cup \dotsm \cup A^n, \forall n \in \mathbb{N}$
        \end{center}
        
        Por tanto podemos notar que de la misma manera seguiremos contando con la cadena vacía $\lambda$. Además al momento de agregar una nueva cadena, ésta seguirá perteneciendo a $A$ y podría seguir siendo la misma.
        
        Si tenemos:
        \begin{center}
            $aaa * aa = aaaaa$,\\ donde\\ $aaa,aa,aaaaa \in A^*$
        \end{center}
        
        También podemos visualizar que la cadena no tendría cambios si alguna de las dos --o las dos-- cuenta con cadena vacía:
        \begin{center}
            $w * \lambda = w$\\
            $\lambda * w = w$\\
            $\lambda * \lambda = \lambda$
        \end{center}
        
        Por lo tanto:
        \begin{center}
            $\boxed{A^* = A^* * A^*}$
        \end{center}
    %
    \item $(A^*)^n = A^*, \forall n \geq 1$.
        \subitem Sea $A = \{a\}$ y $w$ es una cadena de $A$.\\
        
        Tenemos:
        \begin{center}
            $A^* = \{\lambda, a, aa, aaa, \ldots\}$
        \end{center}
        
        Y recordemos también que:
        \begin{center}
            $A^* = A^0 \cup A^1 \cup A^2 \cup \dotsm \cup A^n, \forall n \in \mathbb{N}$
        \end{center}
        
        Tomando en cuenta la propiedad $A^* * A^* = A^*$, que podría reescribirse como $(A^*)^2$, nos indica que la repeticion de eestas concatenaciones no saldría del dominio de $A^*$.
        
        Entonces:
        \begin{center}
            $(A^*)^2 = A^*A^* = A^*$\\
            $(A^*)^3 = (A^*)^2A^* = A^*A^* = A^*$
        \end{center}
        
        Para $n, \forall n \geq 1$:
        \begin{center}
            $(A^*)^n = (A^*)^{n-1} A^* = (A^*)^{n-2}A^*A^* = \dotsm = A^*$
        \end{center}
        
        Por lo tanto, queda demostrado:
        \begin{center}
            $\boxed{A^* = (A^*)^n, \forall n \geq 1}$
        \end{center}
    %
    \item $(A^*)^* = A^*$.
        \subitem Sea $A = \{a\}$. Tenemos que la clausura de Kleene toma todas las repeticiones de $0$ a $n$. Si la aplicamos a $A^*$, seguiría tomando todas las repeticiones. 
        
        Por ello, podríamos tener las cadenas:
        \begin{center}
            $\lambda\lambda\lambda\ldots\lambda = \lambda$\\
            $a * aa * aaa * \dotsm * a...a = aaaaaa...a$
        \end{center}
        Donde:
        \begin{center}
            $aaaaaa,aaa,aa,a,\lambda \in A^*$
        \end{center}
        
        Véase que la operación seguiría sin generar cadenas fuera del dominio de $A^*$.
        
        Por lo tanto:
        \begin{center}
            $\boxed{(A^*)^* = A^*}$
        \end{center}
    %
    \item $A^+ * A^+ \subseteq A^+$.
        \subitem Sea $A = \{a\}$.
        
        Tenemos que:
        \begin{center}
            $A^+ = A^* - \{\lambda\}$
        \end{center}
        
        Por ello, si realizamos la operación, seguiríamos teniendo una cadena perteneciente a la clausura positiva.
        \begin{center}
            $A^+A^+ = A^*A^* - {\lambda}$\\
            y tomando en cuenta la identidad $(A^*)^n = A^*$:\\
            $(A^*)^2 - \{\lambda\} = A^* - \{\lambda\} = A^+$.
        \end{center}
        
        Siendo que:
        \begin{center}
            $A^+ = \{a,aa,aaa,aaaa,\ldots\}$
        \end{center}
        
        Podemos definir que:
        \begin{center}
            $x,y \in A^*,$ donde $x$ y $y$ son cadenas no vacías\\
            $\therefore x \in A^* - \{\lambda\}$; $y \in A^* - \{\lambda\}$\\
            $xy \in A^* - \{\lambda\}$\\
            $\therefore xy \in A^+$
        \end{center}
        
        Por lo tanto:
        \begin{center}
            $\boxed{A^+A^+ \subseteq A^+}$
        \end{center}
    %
    \item $(A^*)^+ = A^*$.
        \subitem Sea $A = \{a\}$. Recordemos que:
        \begin{center}
            $A^* = \bigcup_{i=0}^{\infty}A_i$\\
            $A^+ = \bigcup_{i=1}^{\infty}(A^*)^i$
        \end{center}
        
        Donde notaremos que cada repetición de nuestro $A^*$ tendrá la cadena $\lambda$ aunque el $+$ del paréntesis omita el caso $(A^*)^0$ de nuestro conjunto.\\
        
        Por lo tanto:
        \begin{center}
            $\boxed{(A^*)^+ = A^*}$
        \end{center}
    %
    \newpage
    \item $(A^+)^* = A^*$.
        \subitem Sea $A = \{a\}$. De manera similar a la propiedad $(A^*)^+ = A^*$, definimos lo siguiente:
        
        \begin{center}
            $A^+ = \bigcup_{i=1}^{\infty}A_i$\\
            $A^* = \bigcup_{i=0}^{\infty}(A^+)^i$
        \end{center}
        
        Donde podemos visualizar que nuestra serie ahora toma el caso $(A^+)^0$:
        \begin{center}
            $(A^+)^0 = \{\lambda\}$
        \end{center}
        
        Entonces:
        \begin{center}
            $\lambda \in (A^+)^0$ y $A^* = A^+ \cup \{\lambda\}$\\
        \end{center}
        
        Por lo tanto, podemos decir que:
        \begin{center}
            $\boxed{(A^+)^* = A^*}$
        \end{center}
    %
    \item $(A^+)^+ = A^+$.
        \subitem Sea $A = \{a\}$. De manera similar a la propiedad $(A^*)^* = A^*$, podemos definir:
        
        \begin{center}
            $A^+ = \bigcup_{i=1}^{\infty}A_i$\\
            $(A^+)^+ = \bigcup_{i=1}^{\infty}(A^+)^i$
        \end{center}
        
        Véase que, tanto dentro del paréntesis como fuera, se toma a partir del primer caso, por lo que $\lambda$ nunca llega a estar presente; al igual que en $A^+$.
        
        Por lo tanto, podemos inferir que:
        \begin{center}
            $\boxed{(A^+)^+ = A^+}$
        \end{center}
    %
    \newpage
    \item Si A y B son lenguajes sobre $\Sigma$, entonces $(A \cup B)^* = (A^* B^*)^*$.
        \subitem Sean $A = \{a\}$ y $B = \{b\}$.
        
        Tenemos que:
        \begin{center}
            $A \cup B = \{a,b\}$
        \end{center}
        
        Por lo que entonces:
        \begin{center}
            $(A \cup B)^* = \{a,b,ab,ba,aa,bb,aba,\ldots\}$
        \end{center}
        
        Tomando este último ejemplo que:
        \begin{center}
            $aba \in (A \cup B)^*$
        \end{center}
        
        Para poder igualar, definimos primero que:
        \begin{center}
            $A^* = \{\lambda,a,aa,aaa,aaaa,\ldots\}$\\
            $B^* = \{\lambda,b,bb,bbb,bbbb,\ldots\}$
        \end{center}
        
        Que tiene una forma:
        \begin{center}
            $a^i b^i, \forall i \in \mathbb{N}$
        \end{center}
        
        Por lo que para poder tomar todas las repeticiones tomamos:
        \begin{center}
            $(A^*B^*)^*$\\
            $\therefore aba \in (A^*B^*)^*$
        \end{center}
        
        Por lo tanto, podemos decir que:
        \begin{center}
            $\boxed{(A \cup B)^* = (A^*B^*)^*}$
        \end{center}
    %
\end{enumerate}
    \newpage
    \section{Reflexión o Inverso de un Lenguaje}
        Dado $A$, un lenguaje sobre $\Sigma$, se define $A^R$ de la siguiente forma:
\begin{center}
    $A^R = \{u^R : u \in A\}$,
\end{center}
donde:\\
$u^R$ es la cadena inversa de $u$.\\
$A^R$ se denomina la reflexión o el inverso de $A$.

\textbf{Propiedades:} Sean $A$ y $B$ lenguajes sobre $\Sigma$ (es decir, $A,B \subseteq \Sigma^*$). Demostrar la propiedad 66 y 70.

\begin{enumerate}
    \setcounter{enumi}{65}
    \item $(AB)^R = B^RA^R$.
        \subitem Sea $A = \{a\}$ y $B = \{b\}$.
        
        Tenemos que la concatenación de $A$ y $B$ tendría la forma:
        \begin{center}
            $AB = \{ab\}$
        \end{center}
        
        Cuya forma inversa sería:
        \begin{center}
            $(AB)^R = \{ba\}$
        \end{center}
        
        Si ahora tomamos el lado derecho de la igualdad y lo desarrollamos, tenemos que:
        \begin{center}
            $B^R = \{b\}$\\
            $A^R = \{a\}$
        \end{center}
        
        Cuya concatenación es:
        \begin{center}
            $B^RA^R = \{ba\}$
        \end{center}
        
        Que es igual a la concatenación del lado izquierdo de la igualdad.
        
        Por lo tanto:
        \begin{center}
            $\boxed{(AB)^R = B^R A^R}$
        \end{center}
    %
    \setcounter{enumi}{69}
    \newpage
    \item $(A^*)^R = (A^R)^*$
        \subitem Sea $A = {a}$.
        
        Tenemos:
        \begin{center}
            $A^* = \{\lambda,a,aa,aaa,aaaa,\ldots\}$
        \end{center}
        
        Cuya forma inversa es:
        \begin{center}
            $(A^*)^R = \{a...a,\ldots,aaaa,aaa,aa,a,\lambda\}$
        \end{center}
        
        Si validamos del lado derecho de la igualdad tenemos:
        \begin{center}
            $A^R = \{a\}$
        \end{center}
        
        Cuya clausura de Kleene es:
        \begin{center}
            $(A^R)^* = \{a...a,\ldots,aaaa,aaa,aa,a,\lambda\}$
        \end{center}
        
        Por lo tanto, podemos ver que la igualdad se conserva:
        \begin{center}
            $\boxed{(A^*)^R = (A^R)^*}$
        \end{center}
    %
\end{enumerate}
    \newpage
    \section{Longitud de cadenas}
        Si $w \in \Sigma^*$ y $n,m \in \mathbb{N}$, demuestre que $|w^{n+m}| = |w^n| + |w^m|$.

$R: $

Teniendo que $|w|$ es la magnitud o longitud de una cadena $w$, y tenemos que la potencia de una cadena equivale a la repetición de la misma; es decir, $w^n = \underbrace{w*w*\ldots*w}_{n\text{-veces}}$.

Si continuamos con leyes de los exponentes, podemos notar que de $w^2 = w*w$ y $w^3 = w*w*w$, en su multiplicación se suman sus exponentes:
\begin{center}
    Con $|w| = 1$; $n=2$ y $m=3$:\\
    $|w^2 * w^3| = |\underbrace{\underbrace{w*w}_{2\text{-veces}}*\underbrace{w*w*w}_{3\text{-veces}}}_{5\text{-veces}}|$\\
    $\therefore |w^{n+m}| = |w^n| + |w^m| = 5, n=2, m=3$
\end{center}

Véase que tenemos que $w^n * w^m = w^{n+m}$. Si lo llevamos a términos de magnitud, tenemos que la magnitud de \(w^n\) equivale a la longitud de \(w\) multiplicada por el número de repeticiones. Desarrollando el lado derecho:

\begin{center}
    $|w^n| = |\underbrace{w*\ldots*w}_{n\text{-veces}}| \rightarrow |w^m| = |\underbrace{w*\ldots*w}_{m\text{-veces}}|$\\
    $\therefore |w^n| + |w^m| = |\underbrace{w*\ldots*w}_{n\text{-veces}}| + |\underbrace{w*\ldots*w}_{m\text{-veces}}|$
\end{center}

A partir de esta última premisa y teniendo en cuenta el lado izquierdo de la igualdad ya desarrollado, podemos inferir que:
\begin{center}
    $|w^n| + |w^m| =_{\text{P.Exp.}} n|w| + m|w| =_{\text{distribución}} (n+m)|w|$\\
    y\\
    $|w^{n+m}| =_{\text{P.Exp.}} (n+m)|w| =_{\text{distribución}} n|w| + m|w|$\\
\end{center}

Por lo tanto, queda demostrado:
\begin{center}
    $\boxed{|w^{n+m}| = |w^n| + |w^m|}$
\end{center}
    \newpage
    \section{n-cadenas o n-palabras}
        % ============================================================
% Sección 5: n-cadenas o n-palabras
% ============================================================

\noindent El siguiente teorema expresa el número de $n$-cadenas en función del
número de elementos del alfabeto.

\medskip
\noindent\textbf{Teorema.} Sea $\Sigma$ un alfabeto. Para todo $n \geq 1$ se cumple que
\[
|\Sigma^n| = |\Sigma|^n .
\]

\subsection*{Ejercicio 72}

\noindent\textit{Demuestre el teorema haciendo uso del Principio de Inducción Matemática.}

\medskip
\noindent\textbf{Solución.} Recordemos que $\Sigma^n$ es el conjunto de todas las
cadenas de longitud exactamente $n$ formadas con símbolos de $\Sigma$, es decir,
\[
\Sigma^n = \{\, w \in \Sigma^* : |w| = n \,\}.
\]
Procedemos por inducción sobre $n \geq 1$.

\medskip
\noindent\textbf{1. Base de la inducción ($n = 1$).}
Una cadena de longitud $1$ está formada por un único símbolo, por lo que
\[
\Sigma^1 = \{\, a : a \in \Sigma \,\}.
\]
Cada símbolo $a \in \Sigma$ determina exactamente una cadena de longitud $1$ (la
cadena $a$) y cada cadena de longitud $1$ es un símbolo de $\Sigma$; así, existe una
biyección entre $\Sigma$ y $\Sigma^1$. Por lo tanto
\[
|\Sigma^1| = |\Sigma| = |\Sigma|^1 ,
\]
y la igualdad se cumple para $n = 1$.

\medskip
\noindent\textbf{2. Hipótesis inductiva.}
Supongamos que para un $k \geq 1$ arbitrario se cumple
\[
|\Sigma^k| = |\Sigma|^k .
\]

\medskip
\noindent\textbf{3. Paso inductivo ($n = k+1$).}
Debemos demostrar que $|\Sigma^{k+1}| = |\Sigma|^{k+1}$.

Sea $w \in \Sigma^{k+1}$ una cadena arbitraria. Entonces $w = a_1 a_2 \cdots a_k a_{k+1}$
con $a_i \in \Sigma$ para $1 \leq i \leq k+1$. Podemos expresar $w$ como la
concatenación de dos cadenas más cortas:
\[
w = u\,a, \qquad \text{donde } u = a_1 a_2 \cdots a_k \in \Sigma^k
\ \text{ y } \ a = a_{k+1} \in \Sigma .
\]
Definimos la función
\[
f : \Sigma^k \times \Sigma \longrightarrow \Sigma^{k+1}, \qquad f(u, a) = u\,a .
\]

\begin{itemize}
    \item \textbf{$f$ está bien definida:} si $u \in \Sigma^k$ y $a \in \Sigma$, entonces
    $|ua| = |u| + |a| = k + 1$, por lo que $ua \in \Sigma^{k+1}$.
    \item \textbf{$f$ es inyectiva:} si $f(u,a) = f(u',a')$, entonces $ua = u'a'$. Ambas
    cadenas tienen longitud $k+1$; comparando su último símbolo se obtiene
    $a = a'$, y comparando sus primeros $k$ símbolos se obtiene $u = u'$. Luego
    $(u,a) = (u',a')$.
    \item \textbf{$f$ es suprayectiva:} por la descomposición anterior, toda
    $w = a_1 \cdots a_{k+1} \in \Sigma^{k+1}$ satisface
    $w = f(a_1 \cdots a_k,\ a_{k+1})$.
\end{itemize}

Así, $f$ es una biyección y los conjuntos $\Sigma^{k+1}$ y $\Sigma^k \times \Sigma$
tienen la misma cardinalidad. Por el principio del producto
(una cadena de $\Sigma^{k+1}$ se forma eligiendo primero $u \in \Sigma^k$, lo cual
puede hacerse de $|\Sigma^k|$ maneras, y después $a \in \Sigma$, de $|\Sigma|$
maneras), y usando la hipótesis inductiva:
\[
|\Sigma^{k+1}|
= |\Sigma^k \times \Sigma|
= |\Sigma^k| \cdot |\Sigma|
\overset{\text{H.I.}}{=} |\Sigma|^k \cdot |\Sigma|
= |\Sigma|^{k+1} .
\]

\medskip
\noindent\textbf{Conclusión.} La igualdad se cumple para $n = 1$ y, si se cumple
para $k$, también se cumple para $k+1$. Por el Principio de Inducción Matemática,
\[
|\Sigma^n| = |\Sigma|^n \qquad \text{para todo } n \geq 1. \qquad \blacksquare
\]

\medskip
\noindent\textbf{Observación.} La fórmula también es válida para $n = 0$, pues
$\Sigma^0 = \{\lambda\}$ y $|\Sigma^0| = 1 = |\Sigma|^0$. Por ejemplo, si
$\Sigma = \{0, 1\}$, entonces $|\Sigma^3| = 2^3 = 8$.
    \newpage
    \section{Prefijos y Sufijos de una Cadena}
        % ============================================================
% Sección 6: Prefijos y Sufijos de una Cadena
% ============================================================

\noindent Sean $\Sigma = \{a, b, c, d\}$ y $u = bcbaadb$.

\medskip
\noindent Recordemos que $x$ es \textbf{prefijo} de $u$ si existe $y \in \Sigma^*$ tal
que $u = xy$, y que $y$ es \textbf{sufijo} de $u$ si existe $x \in \Sigma^*$ tal que
$u = xy$. Como $|u| = 7$, la cadena $u$ tiene exactamente un prefijo y un sufijo de
cada longitud $0, 1, \dots, 7$; es decir, $8$ prefijos y $8$ sufijos.

\subsection*{Ejercicio 73}

\noindent\textit{Enumere todos los prefijos de $u$.}

\medskip
\noindent\textbf{Solución.} Tomamos los primeros $i$ símbolos de $u$ para
$i = 0, 1, \dots, 7$:

\begin{center}
    \begin{tabular}{|c|c|c|}
        \hline
        \textbf{Longitud} & \textbf{Prefijo $x$} & \textbf{Resto $y$ ($u = xy$)} \\
        \hline
        0 & $\lambda$  & $bcbaadb$ \\
        1 & $b$        & $cbaadb$ \\
        2 & $bc$       & $baadb$ \\
        3 & $bcb$      & $aadb$ \\
        4 & $bcba$     & $adb$ \\
        5 & $bcbaa$    & $db$ \\
        6 & $bcbaad$   & $b$ \\
        7 & $bcbaadb$  & $\lambda$ \\
        \hline
    \end{tabular}
\end{center}

\[
\text{Pref}(u) = \{\lambda,\ b,\ bc,\ bcb,\ bcba,\ bcbaa,\ bcbaad,\ bcbaadb\}.
\]
Los prefijos propios de $u$ son todos los anteriores excepto $u = bcbaadb$.

\subsection*{Ejercicio 74}

\noindent\textit{Enumere todos los sufijos de $u$.}

\medskip
\noindent\textbf{Solución.} Tomamos los últimos $i$ símbolos de $u$ para
$i = 0, 1, \dots, 7$:

\begin{center}
    \begin{tabular}{|c|c|c|}
        \hline
        \textbf{Longitud} & \textbf{Sufijo $y$} & \textbf{Inicio $x$ ($u = xy$)} \\
        \hline
        0 & $\lambda$  & $bcbaadb$ \\
        1 & $b$        & $bcbaad$ \\
        2 & $db$       & $bcbaa$ \\
        3 & $adb$      & $bcba$ \\
        4 & $aadb$     & $bcb$ \\
        5 & $baadb$    & $bc$ \\
        6 & $cbaadb$   & $b$ \\
        7 & $bcbaadb$  & $\lambda$ \\
        \hline
    \end{tabular}
\end{center}

\[
\text{Suf}(u) = \{\lambda,\ b,\ db,\ adb,\ aadb,\ baadb,\ cbaadb,\ bcbaadb\}.
\]
Los sufijos propios de $u$ son todos los anteriores excepto $u = bcbaadb$.
    \newpage
    \section{Potencias de un Lenguaje}
        % ============================================================
% Sección 7: Potencias de un Lenguaje
% ============================================================

\noindent Dado un lenguaje $A$ sobre $\Sigma$ ($A \subseteq \Sigma^*$) y un número
natural $n \in \mathbb{N}$, se define $A^n$ de la siguiente forma:
\[
A^0 = \{\lambda\}, \qquad
A^n = \underbrace{A \cdot A \cdots A}_{n \text{ veces}}
= \{\, u_1 u_2 \cdots u_n : u_i \in A \text{ para todo } 1 \leq i \leq n \,\}.
\]

\subsection*{Ejercicio 75}

\noindent\textit{Dé una definición recursiva de $A^n$.}

\medskip
\noindent\textbf{Solución.} Usando la concatenación de lenguajes,
$A \cdot B = \{\, uv : u \in A,\ v \in B \,\}$, definimos $A^n$ por recursión sobre
$n \in \mathbb{N}$:
\[
\boxed{
    \begin{aligned}
        &\text{(Base)}      && A^0 = \{\lambda\}, \\
        &\text{(Recursión)} && A^{n+1} = A^n \cdot A, \quad \text{para todo } n \geq 0 .
\end{aligned}}
\]

\medskip
\noindent\textbf{Justificación.} Veamos por inducción sobre $n$ que esta definición
coincide con la definición original:
\[
A^n = \{u_1 \cdots u_n : u_i \in A\}.
\]

\begin{itemize}
    \item Para $n = 0$ ambas definiciones dan $\{\lambda\}$ (la concatenación de cero
    cadenas es la cadena vacía). Además,
    $A^1 = A^0 \cdot A = \{\lambda\} \cdot A = \{\lambda u : u \in A\} = A$.
    \item Si $A^n = \{u_1 \cdots u_n : u_i \in A\}$, entonces
    \[
    A^{n+1} = A^n \cdot A
    = \{\, (u_1 \cdots u_n)\,u_{n+1} : u_i \in A \,\}
    = \{\, u_1 \cdots u_{n+1} : u_i \in A \,\},
    \]
    que es justamente el conjunto de todas las concatenaciones de $n+1$ cadenas
    de $A$.
\end{itemize}

\noindent Como la concatenación es asociativa, también es válida la forma
equivalente $A^{n+1} = A \cdot A^n$.

\medskip
\noindent\textbf{Ejemplo.} Sea $A = \{a, ba\}$. Aplicando la definición recursiva:
\begin{align*}
    A^0 &= \{\lambda\}, \\
    A^1 &= A^0 \cdot A = \{a,\ ba\}, \\
    A^2 &= A^1 \cdot A = \{aa,\ aba,\ baa,\ baba\}, \\
    A^3 &= A^2 \cdot A = \{aaa,\ aaba,\ abaa,\ ababa,\ baaa,\ baaba,\ babaa,\ bababa\}.
\end{align*}
    \newpage
    \section{Ejercicios}
        % ============================================================
% Sección 8: Ejercicios (76 - 93)
% ============================================================

\noindent A lo largo de esta sección usamos las siguientes definiciones recursivas,
válidas para todo $w \in \Sigma^*$ y $a \in \Sigma$:
\[
\begin{aligned}
    &\text{Longitud:}      && |\lambda| = 0, && |wa| = |w| + 1, \\
    &\text{Concatenación:} && u\lambda = u,  && u(wa) = (uw)a .
\end{aligned}
\]
Toda cadena no vacía $v$ puede escribirse de forma única como $v = wa$ con
$w \in \Sigma^*$ y $a \in \Sigma$; esto permite hacer inducción sobre la longitud.

% ------------------------------------------------------------
\subsection*{Ejercicio 76}

\noindent\textit{Sea $\Sigma = \{0, 1\}$. Exhiba $\Sigma^n$ para $n = 2, 3, 4$.}

\medskip
\noindent\textbf{Solución.} $\Sigma^n$ contiene todas las cadenas binarias de
longitud $n$. Las listamos en orden lexicográfico; por el teorema de la
Sección 5 deben ser $2^n$ cadenas.

\begin{align*}
    \Sigma^2 &= \{00,\ 01,\ 10,\ 11\}, && |\Sigma^2| = 2^2 = 4, \\[4pt]
    \Sigma^3 &= \{000,\ 001,\ 010,\ 011,\ 100,\ 101,\ 110,\ 111\}, && |\Sigma^3| = 2^3 = 8, \\[4pt]
    \Sigma^4 &= \{0000,\ 0001,\ 0010,\ 0011,\ 0100,\ 0101,\ 0110,\ 0111, && \\
    &\phantom{= \{} 1000,\ 1001,\ 1010,\ 1011,\ 1100,\ 1101,\ 1110,\ 1111\}, && |\Sigma^4| = 2^4 = 16.
\end{align*}

% ------------------------------------------------------------
\subsection*{Ejercicio 77}

\noindent\textit{Demuestra que $|uv| = |u| + |v|$ para cualesquiera dos cadenas $u$ y $v$.}

\medskip
\noindent\textbf{Demostración.} Fijamos $u \in \Sigma^*$ arbitraria y hacemos
inducción sobre la longitud de $v$.

\medskip
\noindent\textbf{Base ($|v| = 0$, es decir $v = \lambda$).}
\[
|u\lambda| = |u| = |u| + 0 = |u| + |\lambda| .
\]

\noindent\textbf{Hipótesis inductiva.} Supongamos que $|uw| = |u| + |w|$ para toda
cadena $w$ con $|w| = k$.

\medskip
\noindent\textbf{Paso inductivo.} Sea $v$ con $|v| = k + 1$. Entonces $v = wa$ con
$|w| = k$ y $a \in \Sigma$. Usando las definiciones recursivas:
\[
|uv| = |u(wa)| = |(uw)a| = |uw| + 1
\overset{\text{H.I.}}{=} |u| + |w| + 1
= |u| + |wa| = |u| + |v| .
\]
Por el Principio de Inducción Matemática, $|uv| = |u| + |v|$ para todas
$u, v \in \Sigma^*$. $\blacksquare$

% ------------------------------------------------------------
\subsection*{Ejercicio 78}

\noindent\textit{Utiliza inducción para demostrar que $|u^n| = n|u|$ para cualquier
    cadena $u$ y cualquier número natural $n$.}

\medskip
\noindent\textbf{Demostración.} La potencia de una cadena se define recursivamente
como $u^0 = \lambda$ y $u^{n+1} = u^n u$. Hacemos inducción sobre $n$.

\medskip
\noindent\textbf{Base ($n = 0$).} $|u^0| = |\lambda| = 0 = 0 \cdot |u|$.

\medskip
\noindent\textbf{Hipótesis inductiva.} Supongamos que $|u^k| = k|u|$ para algún
$k \geq 0$.

\medskip
\noindent\textbf{Paso inductivo.} Usando el Ejercicio 77:
\[
|u^{k+1}| = |u^k u| \overset{(77)}{=} |u^k| + |u|
\overset{\text{H.I.}}{=} k|u| + |u| = (k+1)|u| .
\]
Por el Principio de Inducción Matemática, $|u^n| = n|u|$ para todo
$n \in \mathbb{N}$. $\blacksquare$

% ------------------------------------------------------------
\subsection*{Ejercicio 79}

\noindent\textit{El reverso de una cadena se define recursivamente como $a^R = a$ y
    $(wa)^R = a\,w^R$ para todo $a \in \Sigma$ y $w \in \Sigma^*$. Utilizando esta
    definición, demuestra que $(uv)^R = v^R u^R$ para todo $u, v \in \Sigma^*$.}

\medskip
\noindent\textbf{Demostración.} Completamos la definición con $\lambda^R = \lambda$
(que es consistente con la dada: tomando $w = \lambda$ se tiene
$a^R = (\lambda a)^R = a\lambda^R = a$). Fijamos $u \in \Sigma^*$ y hacemos
inducción sobre $|v|$.

\medskip
\noindent\textbf{Base ($v = \lambda$).}
\[
(u\lambda)^R = u^R = \lambda\, u^R = \lambda^R u^R .
\]

\noindent\textbf{Hipótesis inductiva.} Supongamos que $(uw)^R = w^R u^R$ para toda
$w$ con $|w| = k$.

\medskip
\noindent\textbf{Paso inductivo.} Sea $v = wa$ con $|w| = k$ y $a \in \Sigma$.
Usando la definición de concatenación, la definición de reverso, la hipótesis
inductiva y la asociatividad de la concatenación:
\[
(uv)^R = (u(wa))^R = ((uw)a)^R
= a\,(uw)^R
\overset{\text{H.I.}}{=} a\,(w^R u^R)
= (a\,w^R)\,u^R
= (wa)^R u^R
= v^R u^R .
\]
Por el Principio de Inducción Matemática, $(uv)^R = v^R u^R$ para todas
$u, v \in \Sigma^*$. $\blacksquare$

% ------------------------------------------------------------
\subsection*{Ejercicio 80}

\noindent\textit{Dado el alfabeto $\Sigma = \{a, b\}$, determina $\Sigma^*$.}

\medskip
\noindent\textbf{Solución.} $\Sigma^*$ es el conjunto de todas las cadenas finitas
sobre $\Sigma$, incluida la cadena vacía:
\[
\Sigma^* = \bigcup_{n \geq 0} \Sigma^n
= \{\lambda,\ a,\ b,\ aa,\ ab,\ ba,\ bb,\ aaa,\ aab,\ aba,\ abb,\ baa,\ \dots\}.
\]
Es un conjunto infinito (numerable), aunque cada una de sus cadenas es finita.

\medskip
\noindent\textit{(a) Da un ejemplo de un lenguaje finito sobre $\Sigma$.}

\medskip
\noindent\textbf{Solución.} Cualquier subconjunto finito de $\Sigma^*$, por ejemplo
\[
L = \{a,\ ab,\ bba\} \subseteq \Sigma^*, \qquad |L| = 3 .
\]

\medskip
\noindent\textit{(b) Dado el lenguaje $L = \{a^n b^n : n \geq 0\}$, determina si las
    cadenas $aabb$, $aaaabbbb$ y $abb$ pertenecen a $L$.}

\medskip
\noindent\textbf{Solución.} Una cadena pertenece a $L$ si y sólo si es un bloque de
$a$'s seguido de un bloque de $b$'s con la misma cantidad de símbolos.
\begin{itemize}
    \item $aabb = a^2 b^2$: $n = 2$. \quad $aabb \in L$.
    \item $aaaabbbb = a^4 b^4$: $n = 4$. \quad $aaaabbbb \in L$.
    \item $abb = a^1 b^2$: tiene una $a$ y dos $b$; como $1 \neq 2$, no es de la
    forma $a^n b^n$. \quad $abb \notin L$.
\end{itemize}

% ------------------------------------------------------------
\subsection*{Ejercicio 81}

\noindent\textit{Dado el lenguaje $L = \{a^n b^n : n \geq 0\}$, calcula: (a) $L^2$,
    (b) $L^R$.}

\medskip
\noindent\textbf{(a)} $L^2 = L \cdot L$ se forma concatenando dos cadenas de $L$,
cuyos exponentes son independientes:
\[
L^2 = \{\, a^n b^n a^m b^m : n, m \geq 0 \,\}
= \{\lambda,\ ab,\ aabb,\ abab,\ aaabbb,\ abaabb,\ aabbab,\ \dots\}.
\]
Nótese que $L \subseteq L^2$ (tomando $m = 0$), pues $\lambda \in L$; pero
$L^2 \neq L$, ya que $abab \in L^2$ y $abab \notin L$.

\medskip
\noindent\textbf{(b)} Para cada $n$, $(a^n b^n)^R = b^n a^n$ (al invertir, el bloque
de $b$'s queda al principio). Entonces
\[
L^R = \{\, b^n a^n : n \geq 0 \,\} = \{\lambda,\ ba,\ bbaa,\ bbbaaa,\ \dots\}.
\]

% ------------------------------------------------------------
\subsection*{Ejercicio 82}

\noindent\textit{Sea $L = \{ab, aa, baa\}$. ¿Cuáles de las siguientes cadenas
    pertenecen a $L^*$?}

\medskip
\noindent\textbf{Solución.} Una cadena pertenece a $L^*$ si puede partirse en
bloques, cada uno igual a $ab$, $aa$ o $baa$. Observemos que la descomposición
está forzada: si el siguiente símbolo es $b$, el bloque sólo puede ser $baa$; si
empieza con $a$, el bloque es $ab$ o $aa$ según el segundo símbolo.

\begin{enumerate}
    \item[(a)] $abaabaaabaa$:
    \[
    abaabaaabaa = (ab)(aa)(baa)(ab)(aa) \in L^* .
    \]
    \item[(b)] $aaaabaaaa$:
    \[
    aaaabaaaa = (aa)(aa)(baa)(aa) \in L^* .
    \]
    \item[(c)] $baaaaabaaaab$: aplicando la descomposición forzada,
    \[
    baaaaabaaaab = (baa)(aa)(ab)(aa)(aa)\ \underbrace{b}_{\notin L^*}
    \]
    Con detalle: el primer bloque debe ser $baa$, quedando $aaabaaaab$;
    luego $aa$, quedando $abaaaab$; luego $ab$, quedando $aaaab$; luego $aa$,
    quedando $aab$; luego $aa$, quedando $b$. La cadena $b$ no es
    concatenación de palabras de $L$ (todas tienen longitud $\geq 2$). Como en
    cada paso la elección era única, no existe ninguna descomposición válida:
    \[
    baaaaabaaaab \notin L^* .
    \]
    \item[(d)] $baaaaabaa$:
    \[
    baaaaabaa = (baa)(aa)(ab)(aa) \in L^* .
    \]
\end{enumerate}

\noindent\textbf{Respuesta:} pertenecen a $L^*$ las cadenas (a), (b) y (d).

% ------------------------------------------------------------
\subsection*{Ejercicio 83}

\noindent\textit{Dado el lenguaje $L = \{a^n b^{n+1} : n \geq 0\}$, ¿es cierto que
    $L^* = L$ para este lenguaje en particular?}

\medskip
\noindent\textbf{Solución.} No. Tenemos
$L = \{b,\ abb,\ aabbb,\ aaabbbb,\ \dots\}$. Basta un contraejemplo:
\begin{itemize}
    \item $\lambda \in L^*$ por definición ($L^0 = \{\lambda\}$), pero
    $\lambda \notin L$, pues toda palabra de $L$ tiene al menos una $b$
    (longitud $2n + 1 \geq 1$).
    \item Incluso sin considerar $\lambda$: $bb = b \cdot b \in L^2 \subseteq L^*$,
    pero $bb \notin L$, porque $bb = a^0 b^2$ y en $L$ el exponente $0$ de $a$
    obliga a exactamente $0 + 1 = 1$ letra $b$.
\end{itemize}
Por lo tanto $L^* \neq L$ (se cumple $L \subsetneq L^*$).

% ------------------------------------------------------------
\subsection*{Ejercicio 84}

\noindent\textit{Dadas las cadenas $u = a^2 b a^3 b^2$ y $v = b a b^2$, calcula lo
    siguiente.}

\medskip
\noindent\textbf{Solución.} Escritas símbolo por símbolo: $u = aabaaabb$ y
$v = babb$.

\begin{enumerate}
    \item[(a)] $uv = aabaaabb \cdot babb = aabaaabbbabb = a^2 b a^3 b^3 a b^2$.
    \item[(b)] $vu = babb \cdot aabaaabb = babbaabaaabb = b a b^2 a^2 b a^3 b^2$.
    \item[(c)] $v^2 = babb \cdot babb = babbbabb = b a b^3 a b^2$.
    \item[(d)] $u^2 = aabaaabb \cdot aabaaabb = aabaaabbaabaaabb
    = a^2 b a^3 b^2 a^2 b a^3 b^2$.
    \item[(e)] $|u| = 2 + 1 + 3 + 2 = 8$.
    \item[(f)] $|v| = 1 + 1 + 2 = 4$.
    \item[(g)] $|uv| = |u| + |v| = 8 + 4 = 12$.
    \item[(h)] $|vu| = |v| + |u| = 4 + 8 = 12$.
    \item[(i)] $|v^2| = 2|v| = 2 \cdot 4 = 8$.
    \item[(j)] $|u^2| = 2|u| = 2 \cdot 8 = 16$.
\end{enumerate}

\noindent Obsérvese que $uv \neq vu$, aunque $|uv| = |vu|$.

% ------------------------------------------------------------
\subsection*{Ejercicio 85}

\noindent\textit{Demuestra que para cualquier par de cadenas $u$ y $v$ se cumple
    que: (a) $|uv| = |u| + |v|$, \ (b) $|uv| = |vu|$.}

\medskip
\noindent\textbf{(a)} Sean $u = a_1 a_2 \cdots a_n$ y $v = b_1 b_2 \cdots b_m$ con
$a_i, b_j \in \Sigma$, de modo que $|u| = n$ y $|v| = m$ (si $n = 0$ o $m = 0$ la
cadena correspondiente es $\lambda$). Por definición de concatenación,
\[
uv = a_1 a_2 \cdots a_n\, b_1 b_2 \cdots b_m ,
\]
que tiene exactamente $n + m$ símbolos. Luego $|uv| = n + m = |u| + |v|$.
(Una demostración formal por inducción sobre $|v|$ se dio en el
Ejercicio 77.) $\blacksquare$

\medskip
\noindent\textbf{(b)} Aplicando el inciso (a) dos veces y la conmutatividad de la
suma en $\mathbb{N}$:
\[
|uv| \overset{(a)}{=} |u| + |v| = |v| + |u| \overset{(a)}{=} |vu| . \qquad \blacksquare
\]
Nótese que la concatenación no es conmutativa ($uv \neq vu$ en general, como en el
Ejercicio 84), pero la longitud sí lo es.

% ------------------------------------------------------------
\subsection*{Ejercicio 86}

\noindent\textit{Dado el alfabeto $A = \{a, b, c\}$, encuentra $L^*$ para:}

\medskip
\noindent\textbf{(a)} $L = \{b^2\}$. Como $L^n = \{(b^2)^n\} = \{b^{2n}\}$,
\[
L^* = \bigcup_{n \geq 0} L^n = \{\, b^{2n} : n \geq 0 \,\}
= \{\lambda,\ bb,\ bbbb,\ bbbbbb,\ \dots\},
\]
es decir, las cadenas formadas sólo por $b$'s con longitud par.

\medskip
\noindent\textbf{(b)} $L = \{a, b\}$. Toda cadena sobre $\{a,b\}$ es concatenación
de sus símbolos, cada uno en $L$; por lo tanto
\[
L^* = \{a, b\}^* = \{\lambda,\ a,\ b,\ aa,\ ab,\ ba,\ bb,\ aaa,\ \dots\},
\]
todas las cadenas sobre $A$ que no contienen $c$.

\medskip
\noindent\textbf{(c)} $L = \{a, b, c^3\}$. Las palabras de $L^*$ se forman
concatenando $a$, $b$ y $ccc$ en cualquier orden y cantidad; por ello las $c$
siempre aparecen en grupos de tres consecutivos:
\[
L^* = \{\, w \in A^* : \text{cada bloque maximal de } c\text{'s en } w
\text{ tiene longitud múltiplo de } 3 \,\}.
\]
Por ejemplo, $\lambda,\ a,\ b,\ ccc,\ abccc,\ cccccca \in L^*$, mientras que
$cc,\ acb \notin L^*$.

% ------------------------------------------------------------
\subsection*{Ejercicio 87}

\noindent\textit{Para el lenguaje $L = \{ab, c\}$ sobre el alfabeto
    $A = \{a, b, c\}$, calcula: (a) $L^3$, (b) $L^{-2}$, (c) $L^0$.}

\medskip
\noindent\textbf{(a)} $L^3 = L \cdot L \cdot L$. Hay $2^3 = 8$ elecciones posibles
(todas producen cadenas distintas):
\[
\begin{aligned}
    L^3 = \{\, & ababab,\ ababc,\ abcab,\ abcc,\\
    & cabab,\ cabc,\ ccab,\ ccc \,\}.
\end{aligned}
\]

\medskip
\noindent\textbf{(b)} La potencia negativa se interpreta con el lenguaje inverso
(reflexión): $L^{-1} = L^R$ y $L^{-n} = (L^{-1})^n = (L^R)^n$. Como
$L^R = \{ba, c\}$:
\[
L^{-2} = (L^R)^2 = \{ba, c\} \cdot \{ba, c\} = \{baba,\ bac,\ cba,\ cc\}.
\]
Comprobación con la propiedad $(A \cdot B)^R = B^R A^R$: $L^2 = \{abab,\ abc,\ cab,\ cc\}$
y $(L^2)^R = \{baba,\ cba,\ bac,\ cc\}$, que coincide con $L^{-2}$.

\medskip
\noindent\textbf{(c)} Por definición, para cualquier lenguaje,
\[
L^0 = \{\lambda\}.
\]

% ------------------------------------------------------------
\subsection*{Ejercicio 88}

\noindent\textit{Dados los lenguajes $L_1 = \{a, ab, a^2\}$ y $L_2 = \{b^2, aba\}$
    sobre el alfabeto $A = \{a, b\}$, determina: (a) $L_1L_2$, (b) $L_2L_2$.}

\medskip
\noindent\textbf{(a)} Concatenamos cada palabra de $L_1$ con cada palabra de $L_2$:

\begin{center}
    \begin{tabular}{|c|c|c|}
        \hline
        $\cdot$ & $bb$ & $aba$ \\
        \hline
        $a$  & $abb$  & $aaba$  \\
        $ab$ & $abbb$ & $ababa$ \\
        $aa$ & $aabb$ & $aaaba$ \\
        \hline
    \end{tabular}
\end{center}
\[
L_1 L_2 = \{abb,\ aaba,\ abbb,\ ababa,\ aabb,\ aaaba\}
= \{ab^2,\ a^2ba,\ ab^3,\ ababa,\ a^2b^2,\ a^3ba\}.
\]

\medskip
\noindent\textbf{(b)}

\begin{center}
    \begin{tabular}{|c|c|c|}
        \hline
        $\cdot$ & $bb$ & $aba$ \\
        \hline
        $bb$  & $bbbb$  & $bbaba$  \\
        $aba$ & $ababb$ & $abaaba$ \\
        \hline
    \end{tabular}
\end{center}
\[
L_2 L_2 = \{bbbb,\ bbaba,\ ababb,\ abaaba\}
= \{b^4,\ b^2aba,\ abab^2,\ aba^2ba\}.
\]

% ------------------------------------------------------------
\subsection*{Ejercicio 89}

\noindent\textit{Dados $u = a^2 b$ y $v = b^3 a b$, encuentra: (a) $uvu$, (b)
    $\lambda u$, $u\lambda$, $u\lambda v$.}

\medskip
\noindent\textbf{Solución.} $u = aab$ y $v = bbbab$.

\medskip
\noindent\textbf{(a)}
\[
uvu = aab \cdot bbbab \cdot aab = aabbbbabaab = a^2 b^4 a b a^2 b,
\qquad |uvu| = 3 + 5 + 3 = 11 .
\]

\noindent\textbf{(b)} Como $\lambda$ es el elemento neutro de la concatenación:
\[
\lambda u = aab = a^2 b, \qquad
u\lambda = aab = a^2 b, \qquad
u\lambda v = uv = aab \cdot bbbab = aabbbbab = a^2 b^4 a b .
\]

% ------------------------------------------------------------
\subsection*{Ejercicio 90}

\noindent\textit{Dado el alfabeto $A = \{a, b, c\}$, determina si $L_1$, $L_2$, $L_3$
    y $L_4$ son lenguajes sobre $A$, donde:}
\[
L_1 = \{a, aa, ab, ac, abc, cab\}, \quad
L_2 = \{aba, aabaa\}, \quad
L_3 = \{\ \}, \quad
L_4 = \{a^i c b^i : i \geq 1\}.
\]

\medskip
\noindent\textbf{Solución.} $L$ es un lenguaje sobre $A$ si y sólo si
$L \subseteq A^*$, es decir, si todas sus palabras usan únicamente símbolos de $A$.
\begin{itemize}
    \item $L_1$: sus palabras sólo usan $a$, $b$ y $c$. \quad \textbf{Sí} es lenguaje
    sobre $A$.
    \item $L_2$: sus palabras sólo usan $a$ y $b$. \quad \textbf{Sí} es lenguaje sobre $A$.
    \item $L_3 = \varnothing$: el conjunto vacío es subconjunto de cualquier conjunto,
    en particular $\varnothing \subseteq A^*$. \quad \textbf{Sí} es lenguaje sobre
    $A$ (el lenguaje vacío; no confundir con $\{\lambda\}$).
    \item $L_4 = \{acb,\ aacbb,\ aaacbbb,\ \dots\}$: sus palabras sólo usan $a$, $b$ y
    $c$. \quad \textbf{Sí} es lenguaje sobre $A$ (infinito).
\end{itemize}

% ------------------------------------------------------------
\subsection*{Ejercicio 91}

\noindent\textit{(a) Dados $L_1 = \{a^i b^j : i > j \geq 1\}$ y
    $L_2 = \{a^i b^j : 1 \leq i < j\}$, encuentra $L_1 \cup L_2$.}

\medskip
\noindent\textbf{Solución.} Ambos lenguajes contienen cadenas $a^i b^j$ con
$i, j \geq 1$; en $L_1$ hay más $a$'s que $b$'s y en $L_2$ hay más $b$'s que $a$'s.
Para $i, j \geq 1$, se cumple $i > j$ o $i < j$ exactamente cuando $i \neq j$. Por lo
tanto
\[
L_1 \cup L_2 = \{\, a^i b^j : i, j \geq 1,\ i \neq j \,\}.
\]
Por ejemplo, $aab,\ abb \in L_1 \cup L_2$, pero $ab,\ aabb \notin L_1 \cup L_2$.

\medskip
\noindent\textit{(b) Dados $L_3 = \{a^i b^i c^j : i, j \geq 1\}$ y
    $L_4 = \{a^i b^j c^j : i, j \geq 1\}$, encuentra $L_3 \cap L_4$.}

\medskip
\noindent\textbf{Solución.} Sea $w \in L_3 \cap L_4$. Entonces
$w = a^i b^i c^j = a^p b^q c^q$ con $i, j, p, q \geq 1$. Como la cadena es un bloque
de $a$'s, seguido de uno de $b$'s y uno de $c$'s, las longitudes de los bloques
coinciden: $i = p$, $i = q$ y $j = q$. Así $i = j$ y $w = a^n b^n c^n$ con
$n \geq 1$. Recíprocamente, toda $a^n b^n c^n$ con $n \geq 1$ está en $L_3$
($i = j = n$) y en $L_4$. Por lo tanto
\[
L_3 \cap L_4 = \{\, a^n b^n c^n : n \geq 1 \,\} = \{abc,\ aabbcc,\ aaabbbccc,\ \dots\}.
\]

% ------------------------------------------------------------
\subsection*{Ejercicio 92}

\noindent\textit{Dados $L_1$ como el lenguaje del inglés y $L_2$ como el lenguaje del
    francés, ¿qué significan (a) $L_1 \cup L_2$ y (b) $L_1 \cap L_2$?}

\medskip
\noindent\textbf{Solución.} Considerando cada idioma como el conjunto de sus
palabras válidas:
\begin{enumerate}
    \item[(a)] $L_1 \cup L_2$ es el conjunto de todas las palabras que son válidas en
    inglés \emph{o} en francés (o en ambos). Por ejemplo, \textit{house},
    \textit{maison} y \textit{table} pertenecen a la unión.
    \item[(b)] $L_1 \cap L_2$ es el conjunto de las palabras que son válidas tanto en
    inglés \emph{como} en francés, por ejemplo palabras prestadas o
    compartidas como \textit{restaurant}, \textit{ballet}, \textit{menu},
    \textit{table} o \textit{question}. En cambio, \textit{house} y
    \textit{maison} no pertenecen a la intersección.
\end{enumerate}

% ------------------------------------------------------------
\subsection*{Ejercicio 93}

\noindent\textit{Sea $A = \{a, b\}$ y $B = \{b, c, d\}$. Definimos los lenguajes}
\[
\begin{aligned}
    L_1 &= \{a^i b^j \mid i \geq 1,\ j \geq 1\}, &
    L_2 &= \{b^i c^j \mid i \geq j \geq 1\}, \\
    L_3 &= \{a^i b^j c^i d^j \mid i \geq 1,\ j \geq 1\}, &
    L_4 &= \{(ad)^i a^j d^j \mid i \geq 2,\ j \geq 1\}.
\end{aligned}
\]
\noindent\textit{Determine si las afirmaciones son verdaderas o falsas.}

\medskip
\noindent\textbf{Solución.} Primero calculamos los alfabetos involucrados:
\[
A \cup B = \{a, b, c, d\}, \quad
A \cap B = \{b\}, \quad
A \oplus B = \{a, c, d\}, \quad
A - B = \{a\}, \quad
B - A = \{c, d\}.
\]
Los símbolos que efectivamente aparecen en cada lenguaje son:
\[
L_1: \{a, b\}, \qquad L_2: \{b, c\}, \qquad L_3: \{a, b, c, d\}, \qquad L_4: \{a, d\}.
\]
Un lenguaje $L$ es lenguaje sobre $\Gamma$ si $L \subseteq \Gamma^*$, es decir, si todos
los símbolos que aparecen en $L$ pertenecen a $\Gamma$. Además, toda palabra de
$L_1$ comienza con $a$ y toda palabra de $L_2$ comienza con $b$, por lo que
$L_1 \cap L_2 = \varnothing$; y el lenguaje vacío es lenguaje sobre cualquier
alfabeto.

\begin{center}
    \renewcommand{\arraystretch}{1.25}
    \begin{tabular}{|c|l|c|l|}
        \hline
        \textbf{Inciso} & \textbf{Afirmación} & \textbf{Valor} & \textbf{Justificación} \\
        \hline
        (a) & $L_1$ sobre $A$ & V & $\{a,b\} \subseteq A$ \\
        (b) & $L_1$ sobre $B$ & F & $a \notin B$ (p.\ ej.\ $ab \in L_1$) \\
        (c) & $L_2$ sobre $A \cup B$ & V & $\{b,c\} \subseteq \{a,b,c,d\}$ \\
        (d) & $L_2$ sobre $A \cap B$ & F & $c \notin \{b\}$ (p.\ ej.\ $bc \in L_2$) \\
        (e) & $L_3$ sobre $A \cup B$ & V & $\{a,b,c,d\} \subseteq A \cup B$ \\
        (f) & $L_3$ sobre $A \cap B$ & F & $a, c, d \notin \{b\}$ \\
        (g) & $L_4$ sobre $A \oplus B$ & V & $\{a,d\} \subseteq \{a,c,d\}$ \\
        (h) & $L_1$ sobre $A - B$ & F & $b \notin \{a\}$ \\
        (i) & $L_1$ sobre $B - A$ & F & $a, b \notin \{c,d\}$ \\
        (j) & $L_1 \cup L_2$ sobre $A$ & F & $c \notin A$ (p.\ ej.\ $bc \in L_2$) \\
        (k) & $L_1 \cup L_2$ sobre $A \cup B$ & V & $\{a,b,c\} \subseteq A \cup B$ \\
        (l) & $L_1 \cup L_2$ sobre $A \cap B$ & F & $a, c \notin \{b\}$ \\
        (m) & $L_1 \cap L_2$ sobre $B$ & V & $L_1 \cap L_2 = \varnothing \subseteq B^*$ \\
        (n) & $L_1 \cap L_2$ sobre $A \cup B$ & V & $\varnothing \subseteq (A \cup B)^*$ \\
        (o) & $L_1 \cap L_2$ sobre $A \cap B$ & V & $\varnothing \subseteq \{b\}^*$ \\
        \hline
    \end{tabular}
\end{center}
\end{document}